% !TeX root = ../navier-stokes-manuscript.tex
% ============================================================
% 2.1 Vortex Stretching
% ============================================================

Axial stretching lengthens one narrow material vortex tube along its axis, contracts its cross-section, draws the rotating fluid inward, and makes the interior spin faster. This spin-up is the positive stretching contribution; the full signed balance with viscosity decides whether the rotation strengthens overall. Both act within the velocity evolution governed by the momentum equation:
\[
  \underbrace{\partial_t u}_{\text{local acceleration}}
  +
  \underbrace{(u\cdot\nabla)u}_{\text{nonlinear transport}}
  =
  \underbrace{-\nabla p}_{\text{pressure}}
  +
  \underbrace{\nu\Delta u}_{\text{viscous diffusion}}.
\]
The tube's net rotation remains undecided, so its geometry and aligned vorticity must be followed together.

\subsection{Vortex Stretching}\label{sub:ThePerfectlyStretchingVortex}

Let \(e\) be the tube axis and \(L(t)\) the length of a small material element within it. Suppose the strain is locally uniform across this element and aligned with \(e\). Writing
\[
  S:=\frac12\bigl(\nabla u+(\nabla u)^{\mathsf T}\bigr),
  \qquad
  Se=\sigma(t)e,
  \qquad
  \sigma(t)>0,
  \qquad
  \frac{\dot L(t)}{L(t)}
  =
  e\cdot Se
  =
  \sigma(t).
\]
The positive strain increases the tube's length at the fractional rate \(\sigma(t)\). Incompressibility forces the same material volume \(\mathcal V\) into a smaller cross-section. Define its effective transverse area by \(A(t):=\mathcal V/L(t)\) and its effective radius by \(r_{\mathrm{eff}}(t):=\sqrt{A(t)/\pi}\). Then
\[
  \nabla\cdot u=0,
  \qquad
  \frac{d}{dt}(A L)=0,
  \qquad
  \frac{\dot A}{A}=-\sigma(t),
  \qquad
  \frac{\dot r_{\mathrm{eff}}}{r_{\mathrm{eff}}}
  =
  -\frac{\sigma(t)}{2}.
\]
The effective radius contracts at half the fractional rate at which the tube lengthens. With \(\omega:=\nabla\times u\), the vorticity carried along the same axis satisfies
\[
  \frac{D\omega}{Dt}
  =
  S\omega+\nu\Delta\omega,
  \qquad
  \frac{D}{Dt}
  :=
  \partial_t+u\cdot\nabla,
\]
and under the alignment \(\omega=|\omega|e\),
\[
  S\omega
  =
  \sigma(t)\omega,
  \qquad
  \frac12\frac{D|\omega|^2}{Dt}
  =
  \sigma(t)|\omega|^2
  +
  \nu\,\omega\cdot\Delta\omega.
\]
The aligned strain supplies positive production, while the viscous term completes the signed decision. On any interval where their sum stays positive, the rotation strengthens as the radius contracts.

The energy identity proved in \Cref{prop:ClassicalKineticEnergyIdentity} and the antisymmetric part of \(\nabla u\) give
\[
  \norm{u(t)}_2
  \leq
  \norm{u_0}_2
  <\infty,
  \qquad
  \left|\frac12\bigl(\nabla u-(\nabla u)^{\mathsf T}\bigr)\right|_{\mathrm F}
  =
  \frac{|\omega|}{\sqrt2}
  \leq
  |\nabla u|_{\mathrm F}.
\]
Finite kinetic energy still permits the rotational part of the gradient to concentrate inside a narrowing region.

\divider{\NavierStokes{} Scaling}

Fix a time \(t_0<T_*\) and call the tube's effective radius then \(\ell:=r_{\mathrm{eff}}(t_0)\). A separate full-field comparison at the finer width \(\lambda^{-1}\ell\), for \(\lambda>1\), uses
\[
  u_\lambda(x,t)
  :=
  \lambda u(\lambda x,\lambda^2t),
  \qquad
  p_\lambda(x,t)
  :=
  \lambda^2p(\lambda x,\lambda^2t).
\]
At time \(\lambda^{-2}t_0\), the corresponding segment in \(u_\lambda\) has effective radius \(\lambda^{-1}\ell\). This is another complete solution at another scale, not a later material state of the opening tube. Rescaling acts on every momentum term:
\[
\begin{aligned}
  \partial_tu_\lambda(x,t)
  &=
  \lambda^3(\partial_tu)(\lambda x,\lambda^2t),
  \\
  (u_\lambda\cdot\nabla)u_\lambda(x,t)
  &=
  \lambda^3\bigl((u\cdot\nabla)u\bigr)(\lambda x,\lambda^2t),
  \\
  \nabla p_\lambda(x,t)
  &=
  \lambda^3(\nabla p)(\lambda x,\lambda^2t),
  \\
  \nu\Delta u_\lambda(x,t)
  &=
  \lambda^3\nu(\Delta u)(\lambda x,\lambda^2t).
\end{aligned}
\]
Rescaling multiplies transport and viscosity by the same cubic factor, so viscosity gains no ground at the finer scale; local acceleration and pressure acquire that factor as well.

\divider{Critical Height}

The momentum balance survives the comparison, while homogeneous Sobolev norms acquire powers that depend on their derivative height. Let \(\Lambda:=(-\Delta)^{1/2}\). Since
\[
  \widehat{u_\lambda}(\xi,t)
  =
  \lambda^{-2}\widehat u(\lambda^{-1}\xi,\lambda^2t),
\]
a change of variables gives
\[
  \norm{u_\lambda(t)}_{\dot H^s}^2
  =
  \lambda^{2s-1}
  \norm{u(\lambda^2t)}_{\dot H^s}^2.
\]
The exponent \(2s-1\) makes kinetic energy lose one power at \(s=0\), makes the first-derivative level gain one at \(s=1\), and vanishes at \(s=\tfrac12\). Set
\[
  H(t)
  :=
  \frac12\norm{u(t)}_{\dot H^{1/2}}^2
  =
  \frac12\norm{\Lambda^{1/2}u(t)}_2^2.
\]
The half-derivative quantity is the critical height for this rescaled family. Writing \(H_\lambda\) for the same quantity evaluated on \(u_\lambda\),
\[
  \norm{u_\lambda(t)}_2^2
  =
  \lambda^{-1}\norm{u(\lambda^2t)}_2^2,
  \qquad
  H_\lambda(t)=H(\lambda^2t),
  \qquad
  \norm{\nabla u_\lambda(t)}_2^2
  =
  \lambda\norm{\nabla u(\lambda^2t)}_2^2.
\]
Rescaling lowers the kinetic energy, raises the first-derivative norm, and leaves the critical height unchanged at corresponding times. Its scale behavior does not decide whether it rises or falls in time.

\divider{Nonlinear Production versus Viscous Dissipation}

At critical height, define the signed nonlinear production by
\[
  \mathcal P_{1/2}[u](t)
  :=
  -\operatorname{Re}
  \pair
    {\Lambda^{1/2}u(t)}
    {\Lambda^{1/2}\bigl((u(t)\cdot\nabla)u(t)\bigr)}_{L^2}.
\]
This production can raise the height. The pressure pairing vanishes because \(\nabla\cdot u=0\), while the Laplacian contributes \(-\nu\norm{\Lambda^{3/2}u}_2^2\). Hence
\[
  H'(t)
  +
  \nu\norm{\Lambda^{3/2}u(t)}_2^2
  =
  \mathcal P_{1/2}[u](t).
\]
Production and positive viscous dissipation decide the sign: \(H'(t)>0\) exactly when production exceeds dissipation. Under the same rescaling,
\[
  \mathcal P_{1/2}[u_\lambda](t)
  =
  \lambda^2\mathcal P_{1/2}[u](\lambda^2t),
  \qquad
  \nu\norm{\Lambda^{3/2}u_\lambda(t)}_2^2
  =
  \lambda^2\nu\norm{\Lambda^{3/2}u(\lambda^2t)}_2^2.
\]
Rescaling multiplies production and dissipation by the same quadratic factor, so neither gains ground. It contracts comparison time by \(t\mapsto\lambda^{-2}t\), but it does not supply the physical interval through which one solution passes from one scale to the next.

\divider{The Heat Clock}

For a linear measure of viscous timing, the heat equation \(v_t=\nu\Delta v\) evolves each Fourier mode by
\[
  \widehat v(\xi,t+\delta t)
  =
  e^{-\nu|\xi|^2\delta t}\widehat v(\xi,t).
\]
The multiplier changes a scale-localized velocity structure of width \(r\), concentrated near \(|\xi|\simeq r^{-1}\), by order one when
\[
  \nu|\xi|^2\delta t\simeq1,
\]
so its linear viscous comparison time is
\[
  \tau_\nu(r)
  :=
  \frac{r^2}{\nu}.
\]
Halving the width quarters the heat clock. More generally,
\[
  \tau_\nu(\lambda^{-1}r)
  =
  \lambda^{-2}\tau_\nu(r).
\]
The linear comparison time contracts by the same factor as time under the full \NavierStokes{} rescaling.

\divider{The Zeno Cascade}

Now sample dyadic widths and their heat times. For
\[
  r_n:=2^{-n}r_0,
  \qquad
  \tau_n:=\frac{r_n^2}{\nu},
\]
the clocks form the finite geometric sum
\[
  \tau_n
  =
  \frac{r_0^2}{\nu}4^{-n},
  \qquad
  \sum_{n=0}^{\infty}\tau_n
  =
  \frac{4r_0^2}{3\nu}
  <
  \infty.
\]
A finite total heat time becomes a possible singular history only if one \NavierStokes{} velocity field exhibits scale-localized structures at widths
\[
  r_0>r_1>r_2>\cdots\longrightarrow0
\]
and at sampled times
\[
  t_0<t_1<t_2<\cdots\uparrow T_*<\infty,
  \qquad
  t_{n+1}-t_n\asymp\tau_n,
\]
with comparison constants independent of \(n\). Under that one-field condition, its remaining time satisfies
\[
  T_*-t_n
  \asymp
  \sum_{k=n}^{\infty}\tau_k
  =
  \frac{4r_n^2}{3\nu}.
\]
The next sampled gap and the total time left then collapse together at order \(r_n^2/\nu\). As those intervals shrink to zero, the unresolved demand falls on that same field: how do its first and successive Eulerian time derivatives behave at one fixed spatial point?
